４）情報通信のシミュレーションの準備課題
1.1000m×1000mのエリアにランダムに端末を20点配置する。
（あまり端末同士が近くならないような配置もできるようにしておく）
2.Random Waypointモデル*により端末を単位時間あたり10m移動させる．
これを1000単位時間繰り返す．
*領域内の目的地点をランダムにさだめ、等速で移動する．目的地に着いたら、新たに目的地を定め移動する．これを繰り返す．斜めの移動がやりにくいなら、最初は縦か横に限定してよい。
3.情報を持つ端末T0をひとつ決め、他の端末が、T0から100m以内にあるとき、T0から情報を受け取る．
4.時間と情報を受け取った端末数をグラフ化する．
　（時間は100,200,...のような一定間隔でよい）
5.情報を受け取った端末が他の端末に情報を渡せるとした場合も同様に調べグラフ化する。
  4.のT0のみの場合と比較する。
6. 情報の受け渡しが100mではなく、50mの場合、150mで同様に調べる。
7.情報を送受信可能な領域をエリアの中央に限定し、情報を受信した端末数と時間をグラフにする．その際
・エリア全体が通信領域の場合と比較する．
・端末数を固定（例えば100）し、通信領域の大きさ(100×100,200×200,...）を可変にして、グラフを比較する．
・通信領域を固定し、端末数を可変にして比較する． 